% ============================================================
%  Chapter 4 — Source Detection: Finding Patient Zero
%  Teaching workbook chapter — builds every concept from scratch
% ============================================================

\chapter{Source Detection: Finding Patient Zero}
\label{chap:source-detection}

\textit{%
  In the previous chapters we learned how epidemic processes spread forward
  through a network: given a starting node, the rules of the SIR model, and a
  contact structure, we can simulate or reason about who gets infected and when.
  This chapter flips the question around completely. We arrive \emph{after the
  fact}, observe who is sick or has recovered, and must reconstruct who started
  it. This is the source detection problem, and it is the central challenge
  this thesis addresses. We build the theory from the very beginning, using
  concrete examples and analogies at every step.
}

% ============================================================
\section{The Detective Analogy}
\label{sec:detective-analogy}
% ============================================================

\begin{analogybox}{The Epidemic Detective}
Imagine you are a detective called to investigate an outbreak. You arrive
\emph{after the fact}: you can see who is currently sick, who has recovered,
and who is still healthy. You have a map of who interacted with whom, and you
know roughly how contagious the disease is. Your job: find \textbf{Patient
Zero} — the single person who introduced the infection.

You walk through the affected community and take stock:

\begin{itemize}
    \item \textbf{Green dots (Susceptible):} people who were never infected.
    \item \textbf{Red dots (Infected):} people who are currently sick.
    \item \textbf{Grey dots (Recovered):} people who were sick and got better.
\end{itemize}

You reason backwards. Patient Zero must have been infected before anyone else,
so they are almost certainly in the recovered group by now. They probably had
many contacts in the early days. The spatial pattern of infection
radiates outward from wherever they are. You look for the ``epicenter.''

This is \emph{exactly} the source detection problem. The detective's reasoning
is formalised as a probabilistic inference task: given the observed snapshot of
states, assign a probability to each person being Patient Zero, and report the
most likely candidate.
\end{analogybox}

The analogy above is not merely illustrative. Every component maps directly
onto a mathematical object:

\begin{center}
\begin{tabular}{ll}
\toprule
\textbf{Detective story} & \textbf{Mathematical object} \\
\midrule
The community of people          & Graph $G = (\mathcal{V}, \mathcal{E})$ \\
Who talked to whom               & Edges $\mathcal{E}$ \\
The snapshot of who is S, I, R   & Observation vector $\mathbf{X}$ \\
Patient Zero                     & Source node $s^*$ \\
``How likely is person $v$ the source?'' & Posterior $P(s = v \mid \mathbf{X})$ \\
The detective's best guess       & MAP estimate $\hat{s} = \arg\max_v P(s=v \mid \mathbf{X})$ \\
\bottomrule
\end{tabular}
\end{center}

The key difficulty — which separates this from a simple puzzle — is that the
SIR spreading process is \textbf{stochastic}. Two runs of the epidemic
starting from the exact same Person Zero can produce different infection
patterns just by chance. This means there is no single ``correct'' answer that
can be deduced with certainty; we must reason probabilistically.

% ============================================================
\section{Formal Problem Statement}
\label{sec:formal-problem}
% ============================================================

Before we can solve a problem, we must state it precisely. This section
introduces the notation used throughout the thesis.

\subsection{The Network}

We model the social contact structure as an undirected graph
$G = (\mathcal{V}, \mathcal{E})$. Here:

\begin{itemize}
    \item $\mathcal{V}$ is the set of \textbf{nodes} (people, animals, computers
    — whatever unit can be infected). We write $N = |\mathcal{V}|$ for the
    total number of nodes. Nodes are labelled $1, 2, \ldots, N$.
    \item $\mathcal{E}$ is the set of \textbf{edges}. An edge $(i, j) \in
    \mathcal{E}$ means nodes $i$ and $j$ were in close enough contact that
    disease transmission is possible.
\end{itemize}

\subsection{The Observation}

After the epidemic has been running for some time, we stop and record the
state of every node. This produces an \textbf{observation vector}:
\[
    \mathbf{X} = (X_1, X_2, \ldots, X_N)
    \quad \text{where each } X_i \in \{S, I, R\}.
\]
$X_i = S$ means node $i$ was never infected (Susceptible). $X_i = I$ means
node $i$ is currently infectious. $X_i = R$ means node $i$ was infected and
has since recovered.

\begin{examplebox}{A Small Observation Vector}
Suppose our network has $N = 6$ nodes and we observe:
\[
    \mathbf{X} = (R,\; R,\; I,\; S,\; R,\; S)
\]
This tells us: nodes 1, 2, 5 have recovered; node 3 is currently sick; nodes
4, 6 were never infected. The infected set (ever infected) is
$\mathcal{I} = \{1, 2, 3, 5\}$. The source must be somewhere in
$\mathcal{I}$ — a node that was never infected cannot have started the
epidemic.
\end{examplebox}

\subsection{The Goal}

\begin{mathbox}{The Source Detection Problem}
\textbf{Given:} A network $G = (\mathcal{V}, \mathcal{E})$, epidemic
parameters $(\beta, \mu)$, and an observation
$\mathbf{X} = (X_1, \ldots, X_N)$ where $X_i \in \{S, I, R\}$.

\textbf{Find:} The node $v^*$ most likely to be the source of the epidemic:
\begin{equation}
    v^* = \arg\max_{v \in \mathcal{V}} \; P\bigl(\text{source} = v \;\big|\; \mathbf{X}\bigr).
    \label{eq:map-estimate}
\end{equation}

Optionally, compute the full \textbf{posterior distribution}
$P(\text{source} = v \mid \mathbf{X})$ for every node $v$, which tells
us how confident we are about each candidate.
\end{mathbox}

The quantity $P(\text{source} = v \mid \mathbf{X})$ is called the
\textbf{posterior probability} of node $v$ being the source, given the
observation. Computing or approximating it is the central technical challenge
of this field.

% ============================================================
\section{Bayes' Theorem: The Mathematical Framework}
\label{sec:bayes}
% ============================================================

The natural tool for this kind of reasoning — updating our belief about an
unknown variable (the source) based on evidence (the observation) — is
\textbf{Bayes' theorem}. If you have not seen it before, here is the full
explanation from first principles.

\subsection{What Bayes' Theorem Says}

Suppose we want to know the probability that event $A$ happened, given that
we observed event $B$. Bayes' theorem states:

\begin{equation}
    P(A \mid B) = \frac{P(B \mid A) \cdot P(A)}{P(B)}
    \label{eq:bayes-general}
\end{equation}

Each term has a name and a meaning:

\begin{conceptbox}{The Four Terms of Bayes' Theorem}
\begin{itemize}
    \item $P(A \mid B)$ — the \textbf{posterior}: the probability of $A$
    \emph{after} we have seen evidence $B$. This is what we want.

    \item $P(B \mid A)$ — the \textbf{likelihood}: how probable is the
    evidence $B$ if $A$ were true? This encodes our model of how $B$ is
    generated.

    \item $P(A)$ — the \textbf{prior}: our belief about $A$ before seeing
    any evidence.

    \item $P(B)$ — the \textbf{evidence}: how probable is the observation
    $B$ overall? It acts as a normalisation constant to make sure the
    posterior sums to 1.
\end{itemize}
\end{conceptbox}

\subsection{Applying Bayes to Source Detection}

In our setting, $A$ is ``node $v$ is the source'' and $B$ is the observation
$\mathbf{X}$. Substituting into Eq.~(\ref{eq:bayes-general}):

\begin{equation}
    \underbrace{P(\text{source}=v \mid \mathbf{X})}_{\text{posterior}}
    \;=\;
    \frac{
        \overbrace{P(\mathbf{X} \mid \text{source}=v)}^{\text{likelihood}}
        \;\times\;
        \overbrace{P(\text{source}=v)}^{\text{prior}}
    }{
        \underbrace{P(\mathbf{X})}_{\text{evidence (normaliser)}}
    }
    \label{eq:bayes-source}
\end{equation}

\textbf{The likelihood} $P(\mathbf{X} \mid \text{source}=v)$ answers:
``If $v$ was the source, how probable is it that we would observe exactly
$\mathbf{X}$?'' This requires us to reason through all the ways the epidemic
could have spread from $v$ to produce $\mathbf{X}$.

\textbf{The prior} $P(\text{source}=v)$ encodes what we believed before
seeing the data. In most cases, we have no reason to suspect one person over
another, so we use a \textbf{uniform prior}:
\[
    P(\text{source}=v) = \frac{1}{N} \quad \text{for all } v \in \mathcal{V}.
\]
This means every node is equally likely a priori. The prior cancels out in
the comparison between candidates, and the MAP estimate simplifies to
\textbf{Maximum Likelihood Estimation (MLE)}:
\[
    v^* = \arg\max_{v \in \mathcal{V}} P(\mathbf{X} \mid \text{source}=v).
\]

\textbf{The evidence} $P(\mathbf{X}) = \sum_{v} P(\mathbf{X} \mid
\text{source}=v) \cdot P(\text{source}=v)$ is just a constant that ensures
all posteriors sum to 1. Since we are only comparing nodes (not computing
absolute probabilities), we can ignore it:
\[
    P(\text{source}=v \mid \mathbf{X})
    \;\propto\;
    P(\mathbf{X} \mid \text{source}=v).
\]

\begin{keyinsight}{The Core Simplification}
With a uniform prior, finding the most likely source is equivalent to finding
the node $v$ from which an epidemic is most likely to produce the
observation $\mathbf{X}$. All our effort therefore goes into computing or
approximating the likelihood $P(\mathbf{X} \mid \text{source}=v)$.
\end{keyinsight}

\subsection{Non-Uniform Priors}

The uniform prior is not always appropriate. If we have additional information
— for example, that the outbreak started at an airport (high-degree hub), or
that there is a known infected traveller — we can use a non-uniform prior to
encode this. In this thesis we restrict ourselves to the uniform prior as the
default, and discuss non-uniform priors only in the context of degree-based
heuristics (Section~\ref{sec:classical-heuristics}).

% ============================================================
\section{Maximum Likelihood Estimation via Simulation}
\label{sec:mle-simulation}
% ============================================================

We have established that the key quantity is the likelihood
$P(\mathbf{X} \mid \text{source}=v)$. But how do we compute it?

\subsection{The Exact Computation is Intractable}

The exact formula for the likelihood requires summing over every possible
way the epidemic could have unfolded from $v$ to produce $\mathbf{X}$.
For a network of $N$ nodes where $k$ got infected, the number of possible
infection sequences (``infection trees'') is proportional to $k^{k-2}$
by Cayley's formula for labelled trees. For $k = 100$, this is already
$100^{98}$ — a number larger than the number of atoms in the observable
universe. Exact computation is hopeless.

\subsection{Monte Carlo Estimation: The Simulation Approach}

Instead, we approximate the likelihood using simulation. The key idea is
elegant:

\begin{conceptbox}{Likelihood by Counting}
If we run $n_{\text{runs}}$ independent SIR simulations starting from
candidate source $v$, and a fraction $f_v$ of them produce an outcome that
``matches'' the observed $\mathbf{X}$, then:
\[
    P(\mathbf{X} \mid \text{source}=v) \approx f_v.
\]
The more simulations we run, the more accurate the estimate.
\end{conceptbox}

This is the \textbf{Monte Carlo maximum likelihood estimator}. It powers the
Soft Margin Estimator (SME), one of the baseline methods in this thesis.

\subsection{A Worked Example: MLE on a 4-Node Network}

\begin{examplebox}{MLE by Simulation on a Small Network}
Consider a network with $N = 4$ nodes arranged as a path:
$1 - 2 - 3 - 4$. The observed state is:
\[
    \mathbf{X}_{\text{obs}} = (R,\; R,\; I,\; S)
\]
Nodes 1 and 2 have recovered, node 3 is still infected, node 4 was never
infected. We want to estimate the likelihood of each candidate source.

We run $n_{\text{runs}} = 5$ simulations from each candidate and record the
outcome (simplified to show just whether each node is in $\{S, I, R\}$):

\bigskip
\textbf{Simulations from $v = 1$:}
\begin{center}
\begin{tabular}{cccccc}
\toprule
\textbf{Run} & \textbf{Node 1} & \textbf{Node 2} & \textbf{Node 3} & \textbf{Node 4} & \textbf{Match?} \\
\midrule
1 & R & R & I & S & \checkmark match \\
2 & R & S & S & S & $\times$ \\
3 & R & R & R & I & $\times$ \\
4 & R & R & I & S & \checkmark match \\
5 & R & R & S & S & $\times$ \\
\bottomrule
\end{tabular}
\end{center}
2 out of 5 runs match $\mathbf{X}_{\text{obs}}$.
Estimated likelihood: $\hat{P}(\mathbf{X}_{\text{obs}} \mid v=1) = 2/5 = 0.40$.

\bigskip
\textbf{Simulations from $v = 4$:}
Node 4 ends up in state $S$ in $\mathbf{X}_{\text{obs}}$. A source node
\emph{always} starts as Infected and transitions to Recovered — it can never
be Susceptible at observation time. So \emph{every} simulation from $v = 4$
will have node 4 in state $R$ or $I$, never $S$. No simulation can match.
Estimated likelihood: $\hat{P}(\mathbf{X}_{\text{obs}} \mid v=4) = 0/5 = 0$.

\bigskip
This illustrates the important \textbf{candidate filtering} rule we formalise
later: nodes that are Susceptible in $\mathbf{X}_{\text{obs}}$ can be
immediately eliminated as candidates.
\end{examplebox}

\subsection{The Computational Cost of Brute-Force MLE}

This simulation approach is correct and unbiased, but its cost grows rapidly
with network size. We examine the scaling in detail in
Section~\ref{sec:computational-challenge}. For now, note that for each of the
$N$ candidate sources we must run $n_{\text{runs}}$ full epidemic simulations.
If each simulation takes time $t_{\text{sim}}$, the total cost is:
\[
    \text{Total time} = N \times n_{\text{runs}} \times t_{\text{sim}}.
\]
This is the motivation for every faster method in the remainder of the thesis.

% ============================================================
\section{The Soft Matching Problem}
\label{sec:soft-matching}
% ============================================================

The MLE approach above used \textbf{exact matching}: a simulation counts as a
``match'' only if every single node's state agrees with $\mathbf{X}_{\text{obs}}$.
This is a natural choice but can be too strict in practice.

\subsection{Why Exact Matching Fails}

Consider an epidemic on a network with $N = 1000$ nodes. A single simulated
run produces an outcome that agrees with the observation on 998 out of 1000
nodes — extremely close! — but the two disagreeing nodes push it into the
``no match'' category, contributing nothing to the likelihood estimate.

More fundamentally:

\begin{itemize}
    \item \textbf{Observation noise:} In practice, disease states are
    self-reported or imperfectly tested. Some observed states may be wrong.
    Requiring exact agreement with a noisy observation is counter-productive.
    \item \textbf{Stochastic near-misses:} A simulation from the true source
    has a high probability of being \emph{close} to the observation, but exact
    agreement may be rare even when the source is correct.
    \item \textbf{Variance reduction:} By quantifying the \emph{degree} of
    agreement rather than demanding binary match, we extract more signal from
    each simulation run.
\end{itemize}

\subsection{Soft Matching with Overlap Metrics}

Instead of a binary match, we measure how similar the simulated outcome is to
the observation. A natural measure is the \textbf{Jaccard similarity} between
the two infected sets.

Let $\mathcal{I}_{\text{obs}}$ be the set of nodes that are infected or
recovered in the observation, and let $\mathcal{I}_{\text{sim}}$ be the
corresponding set from a single simulation run. The Jaccard similarity is:

\begin{mathbox}{Jaccard Similarity}
\begin{equation}
    J(\mathcal{I}_{\text{obs}},\, \mathcal{I}_{\text{sim}})
    = \frac{|\mathcal{I}_{\text{obs}} \cap \mathcal{I}_{\text{sim}}|}
           {|\mathcal{I}_{\text{obs}} \cup \mathcal{I}_{\text{sim}}|}
    \label{eq:jaccard}
\end{equation}
This equals 1 if the two sets are identical, 0 if they share no elements,
and takes values in $(0, 1)$ otherwise. It penalises both false positives
(nodes in $\mathcal{I}_{\text{sim}}$ not in $\mathcal{I}_{\text{obs}}$) and
false negatives (nodes in $\mathcal{I}_{\text{obs}}$ not in
$\mathcal{I}_{\text{sim}}$).
\end{mathbox}

\begin{examplebox}{Jaccard Similarity — Numerical Computation}
Let the observation infect nodes $\mathcal{I}_{\text{obs}} = \{1, 2, 3\}$
and a simulation produce $\mathcal{I}_{\text{sim}} = \{2, 3, 4\}$.

\begin{itemize}
    \item Intersection: $\mathcal{I}_{\text{obs}} \cap \mathcal{I}_{\text{sim}}
    = \{2, 3\}$, so $|\cdot| = 2$.
    \item Union: $\mathcal{I}_{\text{obs}} \cup \mathcal{I}_{\text{sim}}
    = \{1, 2, 3, 4\}$, so $|\cdot| = 4$.
    \item Jaccard similarity: $J = 2 / 4 = 0.5$.
\end{itemize}

This means 50\% overlap between the two infected sets. Contrast with a
simulation giving $\mathcal{I}_{\text{sim}} = \{1, 2, 3\}$: then
$J = 3/3 = 1.0$ (perfect match), or
$\mathcal{I}_{\text{sim}} = \{5, 6, 7\}$: then $J = 0/6 = 0$ (no overlap).
\end{examplebox}

\subsection{The Soft Margin Estimator (SME)}

The Soft Margin Estimator, which serves as a strong baseline in this
thesis, uses soft matching as a proxy likelihood:

\[
    \hat{P}_{\text{soft}}(\mathbf{X}_{\text{obs}} \mid \text{source}=v)
    \;\approx\; \frac{1}{n_{\text{runs}}}
    \sum_{r=1}^{n_{\text{runs}}} J\!\bigl(\mathcal{I}_{\text{obs}},\;
    \mathcal{I}^{(r)}_v\bigr)
\]

where $\mathcal{I}^{(r)}_v$ is the infected set in simulation run $r$ from
candidate source $v$. This averages the Jaccard similarity over all runs,
giving a continuous score in $[0,1]$ that varies smoothly with source
quality.

\begin{keyinsight}{Soft vs.\ Hard Matching}
Exact matching is the theoretically correct likelihood estimator (under a
perfect observation model), but it has very high variance: most simulations
produce zero contribution. Soft matching introduces a small bias (it is no
longer an exact likelihood) but drastically reduces variance, making it
practical even with a modest number of simulation runs.
\end{keyinsight}

% ============================================================
\section{Challenges in Source Detection}
\label{sec:challenges}
% ============================================================

Armed with the Bayesian framework and the simulation-based estimator, we can
now examine why source detection is hard in practice. Each challenge below
is a fundamental reason why simple methods fail, and each motivates a design
choice in the GNN architectures developed later in this thesis.

\subsection{Challenge 1: Partial Observation}

\begin{conceptbox}{Challenge: We Cannot See Everyone}
In the idealized setting, we observe the state of \emph{every} node in the
network. In practice, this is rarely possible:
\begin{itemize}
    \item In a disease outbreak, only symptomatic patients seek care. Mild or
    asymptomatic cases are invisible.
    \item In a computer network intrusion, only monitored nodes generate logs.
    \item In a social network rumour spread, only users who engage publicly
    are visible.
\end{itemize}
\textbf{Why it makes things harder:} Suppose the true source $s^*$ infected
an intermediate node $u$ early on, and $u$ infected three others. If node
$u$ is unobserved, we lose the connection between $s^*$ and the observed
infected nodes. Another candidate might explain the observed infected nodes
just as well, even if it required an improbable transmission path.

\textbf{Analogy:} The detective's crime scene has missing evidence — some
crucial footprints have been walked over. The reconstruction becomes much
less certain.
\end{conceptbox}

\begin{examplebox}{Partial Observation Makes Multiple Sources Plausible}
Network: $A - B - C - D - E$. True source: $A$.

\textbf{Full observation:} $\mathbf{X} = (R, R, R, I, S)$.
Node $A$ is the only candidate geometrically consistent with a wave of
infection spreading left-to-right.

\textbf{Partial observation} (only nodes $C$, $D$, $E$ visible):
$\mathbf{X}_{\text{partial}} = (?, ?, R, I, S)$.
Now both $A$ and $B$ appear equally plausible as sources from the visible
evidence, because $A$ and $B$ are both ``upstream'' of $C$.
\end{examplebox}

\subsection{Challenge 2: Late Observation}

\begin{conceptbox}{Challenge: The Signal Fades with Time}
If we observe the epidemic long after it started, the infected region has
spread far from the source. The source is now deeply embedded in a large
recovered region, and the boundary of the infected set carries less
information about the origin.

\textbf{Analogy:} Consider a drop of ink in a glass of water. Just after the
drop falls, the dark spot pinpoints where it landed. One minute later, the
ink has diffused evenly and there is no way to tell where it entered.

\textbf{Why it matters for our methods:} Most classical methods (like the
Jordan center) work well early, when the infected set is small and
geographically concentrated. As $T$ grows, the infection expands nearly
uniformly in all directions and any node deep inside the recovered region is
equally plausible. The posterior becomes flat.

\textbf{How GNNs help:} A GNN trained on epidemic trajectories can learn
subtler signatures of the source — not just ``is this node central?'' but
``does this node's local neighbourhood pattern match what the source's
neighbourhood typically looks like?''
\end{conceptbox}

\subsection{Challenge 3: Network Heterogeneity}

Real networks are not uniform. A social network has a few people with
hundreds of contacts (hubs) and many people with only a handful. Consider
two scenarios:

\begin{examplebox}{Hub vs.\ Leaf as Source}
Suppose Node $H$ is a hub with degree 50, and Node $L$ is a leaf with
degree 1. Both are infected.

\textbf{If $H$ is the true source:} The epidemic spreads rapidly outward from
$H$ in all directions. The resulting infected set looks like a large,
roughly symmetric ball centred at $H$. This pattern is \emph{distinctive} —
it is hard to reproduce from a leaf node. $H$ is easy to detect.

\textbf{If $L$ is the true source:} The epidemic enters through $L$'s single
edge to its neighbour, and then fans out from there. The infected set looks
as if it started from $L$'s one neighbour, not $L$ itself. The posterior
over $L$ vs.\ its neighbour is nearly flat.

\textbf{Consequence:} Detection accuracy varies systematically with the
source's degree. High-degree sources are easier to detect; low-degree
sources blend in with their neighbours.
\end{examplebox}

\subsection{Challenge 4: Temporal Ambiguity}

In a static network, two different source nodes may produce
statistically identical observation distributions. This is especially true
on homogeneous networks like regular grids or random regular graphs, where
every node has the same local structure.

\begin{examplebox}{Symmetric Networks Kill Identifiability}
Consider a 4-cycle: $1 - 2 - 3 - 4 - 1$. Every node is structurally
identical (same degree, same local neighbourhood). If the observation is
$(R, I, S, R)$, both node 1 and node 4 are equally plausible sources by
symmetry — the network has a reflection symmetry that maps each scenario
onto the other. No method can distinguish them better than 50:50 on such
a symmetric graph.

\textbf{Temporal networks break this symmetry.} If the actual timing of
contacts is recorded, the two scenarios may differ: perhaps the edge
$(4, 1)$ was active at $t=1$ and the edge $(1, 2)$ at $t=3$, making it
causally possible for 4 to infect 1 first, but not the reverse. Temporal
information restores identifiability.
\end{examplebox}

\begin{keyinsight}{Why Temporal Information Matters}
Adding contact timestamps to the network model can transform an ambiguous
inference problem into a near-certain one. This is the central motivation
for temporal GNN architectures: they explicitly model the temporal
ordering of contacts, dramatically reducing the set of causally plausible
source candidates. Chapter~2 introduced causal walks; later chapters show
how GNNs learn to exploit them.
\end{keyinsight}

% ============================================================
\section{Evaluation Metrics}
\label{sec:evaluation-metrics}
% ============================================================

How do we know whether a source detection method is good? We need metrics
that capture performance across many test cases, handle varying network
sizes, and measure whether the method's top predictions include the true
source. This section defines every metric used in the codebase, with full
worked examples.

\subsection{Setup: Scores and Ranks}

After running a source detection method on a test case, we obtain a
\textbf{score vector} $\mathbf{p} = (p_1, p_2, \ldots, p_N)$ where $p_v$
is the method's predicted probability (or proxy score) for node $v$ being
the source. Higher $p_v$ means the method considers $v$ more likely.

We then \textbf{rank} the nodes from most likely to least likely. If we sort
in descending order of $p_v$, the node with the highest score gets rank 1,
the second highest gets rank 2, and so on.

\begin{examplebox}{From Scores to Ranks}
Suppose $N = 5$ and the method outputs scores:
\[
    \mathbf{p} = (0.05,\; 0.60,\; 0.25,\; 0.08,\; 0.02)
\]
Sorted in descending order: node 2 ($0.60$), node 3 ($0.25$), node 4
($0.08$), node 1 ($0.05$), node 5 ($0.02$).

Ranks: node 1 gets rank 4, node 2 gets rank 1, node 3 gets rank 2, node 4
gets rank 3, node 5 gets rank 5.

If the true source is node 3, then the true source has \textbf{rank 2} in
this prediction.
\end{examplebox}

\subsection{Rank Score}

The \textbf{rank score} measures where the true source appears in the sorted
prediction list, normalised by the network size.

\begin{mathbox}{Rank Score}
Let $r_{s^*}$ be the rank of the true source $s^*$ in the method's sorted
predictions (rank 1 = most likely). The rank score for a single test case is:
\begin{equation}
    \text{RankScore} = \frac{r_{s^*}}{N}
    \label{eq:rank-score}
\end{equation}
\begin{itemize}
    \item \textbf{Best possible:} $r_{s^*} = 1 \Rightarrow
    \text{RankScore} = 1/N \approx 0$ (true source ranked first).
    \item \textbf{Worst possible:} $r_{s^*} = N \Rightarrow
    \text{RankScore} = 1$ (true source ranked last).
    \item \textbf{Random baseline:} Expected rank of a random guess is
    $(N+1)/2$, giving expected RankScore $\approx 0.5$.
\end{itemize}
\textbf{Lower is better.}
\end{mathbox}

\begin{examplebox}{Rank Score — Worked Computation}
Suppose $N = 100$ and the true source ends up ranked 5th out of 100.
\[
    \text{RankScore} = \frac{5}{100} = 0.05.
\]
This is excellent: the method placed the true source in the top 5\% of
candidates. Compare to a random method which would give
$\text{RankScore} \approx 0.5$ on average.

In the codebase (\texttt{eval/scores.py}), this is computed slightly
differently to allow for a configurable offset parameter:
\begin{center}
\texttt{rank\_score(ranks, sel, offset=0)} $=$ \texttt{mean((1 + offset) / (ranks[sel] + offset))}
\end{center}
With \texttt{offset=0} this gives $1/r_{s^*}$, the \textbf{reciprocal rank}.
Note: reciprocal rank is also common in information retrieval — higher is
better (range $[1/N, 1]$). Whether you use rank/$N$ (lower is better) or
$1/\text{rank}$ (higher is better) is a matter of convention; the codebase
uses the reciprocal rank formulation.
\end{examplebox}

\begin{codewalk}{Rank Score in \texttt{eval/scores.py}}
\begin{lstlisting}[language=Python]
def rank_score(ranks, sel, offset=0):
    """
    ranks : array of shape (n_sources * n_runs,)
            rank of the true source in each test case.
    sel   : boolean mask selecting which test cases to include
            (e.g. only cases with maximal_outbreak==True).
    offset: optional smoothing term (default 0).

    Returns the mean reciprocal rank over selected cases.
    With offset=0: mean(1 / ranks[sel])
    """
    return np.mean((1 + offset) / (ranks[sel] + offset))
\end{lstlisting}

The \texttt{sel} mask is crucial: it restricts evaluation to the subset of
test cases we consider valid (see Section~\ref{sec:candidate-filtering}).
The \texttt{ranks} array holds the rank of the true source node for each
test case; a rank of 1 means the true source was the top prediction.
\end{codewalk}

\subsection{Top-$k$ Score}

The \textbf{top-$k$ score} is a binary metric that asks: ``Is the true
source within the top $k$ predictions?''

\begin{mathbox}{Top-$k$ Score}
\begin{equation}
    \text{Top-}k\text{Score} = \mathbf{1}\bigl[r_{s^*} \leq k\bigr]
    \in \{0, 1\}
    \label{eq:topk-score}
\end{equation}
where $\mathbf{1}[\cdot]$ is the indicator function (1 if true, 0 if false).

Over many test cases, the \textbf{mean top-$k$ score} is the fraction of
test cases where the true source appeared in the top $k$ predictions. This
is analogous to ``recall at $k$'' in information retrieval.

\textbf{Higher is better.} A perfect method has mean top-$k$ score of 1.0.
A random method has mean top-$k$ score of $k/N$.
\end{mathbox}

\begin{examplebox}{Top-$k$ Score — Worked Example}
Suppose $N = 100$, $k = 5$, and we have 4 test cases:

\begin{center}
\begin{tabular}{ccc}
\toprule
\textbf{Test case} & \textbf{True source rank} & \textbf{Top-5 score} \\
\midrule
1 & 3  & 1 \\
2 & 12 & 0 \\
3 & 1  & 1 \\
4 & 5  & 1 \\
\bottomrule
\end{tabular}
\end{center}

Mean top-5 score $= (1 + 0 + 1 + 1)/4 = 0.75$.

The method places the true source in the top 5 in 3 out of 4 cases. Compare
to random: expected score $= 5/100 = 0.05$.
\end{examplebox}

\begin{codewalk}{Top-$k$ Score in \texttt{eval/scores.py}}
\begin{lstlisting}[language=Python]
def top_k_score(ranks, sel, k=5):
    """
    Returns the fraction of selected test cases where
    the true source is ranked <= k (i.e., in the top-k).

    ranks[sel] <= k  produces a boolean array;
    np.mean() over booleans gives the fraction of True values.
    """
    return np.mean(ranks[sel] <= k)
\end{lstlisting}

Notice the elegant simplicity: \texttt{ranks[sel] <= k} is a boolean array,
and \texttt{np.mean} over booleans equals the fraction that are \texttt{True}.
\end{codewalk}

\subsection{Mean Rank Score}

In practice, we evaluate a method on many test cases (many different epidemic
realisations, each with a known true source) and report the
\textbf{mean rank score} — the average of the rank score across all test
cases in the evaluation set.

\begin{mathbox}{Mean Rank Score}
Given $M$ test cases with true source ranks $r_1, r_2, \ldots, r_M$:
\begin{equation}
    \overline{\text{MRR}}
    = \frac{1}{M} \sum_{m=1}^{M} \frac{1}{r_m}
    \label{eq:mean-rr}
\end{equation}
(using the reciprocal rank convention from the codebase). A method achieving
$\overline{\text{MRR}} = 1.0$ always ranks the true source first.
A random baseline achieves $\overline{\text{MRR}} = \frac{2}{N+1} \approx
2/N$ for large $N$.
\end{mathbox}

\subsection{Additional Metrics in the Codebase}

The codebase (\texttt{eval/scores.py}) also implements several additional
metrics for richer evaluation:

\begin{conceptbox}{Additional Evaluation Metrics}
\textbf{Normalised Brier Score.} Measures the squared error between the
predicted probability vector $\mathbf{p}$ and the one-hot true label vector.
It penalises both wrong predictions and miscalibrated confidence. Lower is
better. Normalised by the expected Brier score of a uniform predictor
($2/N$), so a score of 1.0 means the method is no better than uniform.

\textbf{Normalised Entropy.} Measures how spread out (uncertain) the
method's predicted distribution is. A method that concentrates probability
on few nodes has low entropy (high confidence). Normalised by the entropy of
the uniform distribution ($\log N$), so a score of 1.0 means maximum
uncertainty. Useful for quantifying confidence rather than accuracy.

\textbf{Credible Set Coverage.} Computes the smallest set of nodes that
together contain at least probability mass $p$ (e.g., $p = 0.9$), and checks
whether the true source is in this ``credible set.'' A small credible set
that contains the true source is ideal. This metric penalises methods that
need to hedge across many nodes.
\end{conceptbox}

% ============================================================
\section{Candidate Source Filtering}
\label{sec:candidate-filtering}
% ============================================================

Before ranking all $N$ nodes, we can immediately eliminate many candidates
using logical constraints. This section explains the two main filters used
in the codebase.

\subsection{The ``Possible'' Filter}

\begin{conceptbox}{Who Can Be Patient Zero?}
A node can only be the epidemic source if it was \emph{at some point
infected}. A node that is Susceptible at observation time — meaning it was
never infected during the entire epidemic — cannot possibly be the source,
because the source is always the first infected node.

\textbf{Formally:} Node $v$ is a \textbf{possible} candidate source in a
specific epidemic realisation if and only if $v \in \mathcal{I}_{\text{obs}}$,
i.e., $X_v \in \{I, R\}$ at observation time.
\end{conceptbox}

In the codebase, this is encoded in the \texttt{possible} array loaded by
\texttt{setup/data\_loader.py}:
\begin{itemize}
    \item Shape: \texttt{(n\_nodes, n\_runs, n\_nodes)}. For each true source
    $s$ (first axis), each run (second axis), and each candidate node $v$
    (third axis), \texttt{possible[s, run, v] = 1} if $v$ was infected at
    some point in this run.
    \item A candidate node $v$ that is never infected has
    \texttt{possible[s, run, v] = 0} and should be masked out before ranking.
\end{itemize}

\begin{codewalk}{Candidate Filtering in \texttt{setup/data\_loader.py}}
\begin{lstlisting}[language=Python]
# possible[s, run, v] == 1 if node v was infected in run 'run'
# when the true source was 's'
possible = np.fromfile(
    f"data/{tsir_run_id}/possible_sources.bin",
    dtype=np.int8
).reshape(n_nodes, n_runs, n_nodes)

# lik_possible is used to mask out impossible candidates from
# log-likelihoods: 0 for possible nodes, +inf for impossible
# (adding +inf to a log-likelihood makes it -inf after negation,
#  ensuring impossible nodes are never ranked first)
lik_possible = np.where(possible == 1, 0, np.inf)
\end{lstlisting}

The \texttt{lik\_possible} trick is elegant: adding this to a log-likelihood
score sets impossible candidates to $-\infty$ log-probability, ensuring they
appear last in the ranking regardless of the raw model score.
\end{codewalk}

\begin{examplebox}{Effect of the Possible Filter}
Suppose $N = 10$ and the observation has nodes $\{3, 5, 7, 9\}$ infected or
recovered. The remaining 6 nodes are Susceptible.

\textbf{Without filtering:} We rank all 10 nodes. Even if the model correctly
scores node 3 highest, a bad method might accidentally score a Susceptible
node higher, wasting prediction budget.

\textbf{With filtering:} Only nodes $\{3, 5, 7, 9\}$ are candidates. The
effective problem size drops from 10 to 4. All scores for Susceptible nodes
are set to $-\infty$. The method cannot accidentally predict an impossible
candidate.

\textbf{Quantitative impact:} If the true source is uniformly distributed
among the $k$ possible candidates, the random baseline rank score improves
from $\approx N/2$ to $\approx k/2$. In a typical epidemic where 40\% of
nodes are infected, this cuts the effective search space by 60\%.
\end{examplebox}

\subsection{The ``Maximal Outbreak'' Filter}

Not all epidemic realisations are equally informative for evaluation. If the
epidemic barely spreads (e.g., only 2 nodes get infected in a 1000-node
network), the observation carries almost no information about the source:
the source and its immediate neighbour both look identical. Including such
trivial cases would artificially inflate performance metrics.

\begin{conceptbox}{The Maximal Outbreak Filter}
The \textbf{maximal outbreak} filter restricts evaluation to epidemic
realisations where the epidemic spread significantly — typically defined as
infecting at least some threshold fraction of the network (e.g., $\geq 10\%$
of nodes).

\textbf{Why filter?} Trivial outbreaks do not test the method's ability to
distinguish between candidates; they are essentially random. Including them
makes all methods look equally bad (or equally good) and obscures the
differences between methods on informative cases.

In the codebase, \texttt{maximal\_outbreak} is a boolean mask of shape
\texttt{(n\_sources, n\_nodes)} that is \texttt{True} for
source--realisation pairs where the epidemic reached a significant fraction
of the network. The scoring functions receive a \texttt{sel} parameter
(selection mask) that combines the maximal outbreak filter with any other
desired restrictions.
\end{conceptbox}

\begin{codewalk}{Maximal Outbreak Mask in \texttt{setup/data\_loader.py}}
\begin{lstlisting}[language=Python]
# maximal_outbreak[s, v] = 1 if the epidemic starting from source s
# infected a significant fraction of the network in a maximal run.
# Shape: (n_nodes, n_nodes)
maximal_outbreak = (
    1 - np.fromfile(
        f"data/{tsir_run_id}/maximal_outbreak_S.bin",
        dtype=np.int8
    ).reshape(n_nodes, n_nodes)
)
# The file stores whether node v is Susceptible in the maximal run;
# subtracting from 1 flips it to "was node v infected?".
# maximal_outbreak[s, v] = 1 means node v was infected in the
# maximal run from source s, i.e., the epidemic was large enough.
\end{lstlisting}
\end{codewalk}

The following pseudocode shows how the two filters combine in practice:

\begin{codewalk}{Applying Both Filters Before Scoring}
\begin{lstlisting}[language=Python]
# Suppose we have:
#   scores: array (n_sources, n_runs, n_nodes) -- raw model output
#   possible: (n_sources, n_runs, n_nodes)     -- possible mask
#   maximal_outbreak: (n_sources, n_nodes)     -- outbreak mask

# Step 1: mask out impossible candidates per run
# (set their scores to -inf so they rank last)
scores_filtered = scores - data.lik_possible   # lik_possible = 0 or +inf

# Step 2: rank nodes within each run (descending by score)
ranks = np.argsort(-scores_filtered, axis=-1)  # shape: (n_sources, n_runs, n_nodes)
# find rank of true source in each run
true_source_ranks = ...   # find where true source sits in ranking

# Step 3: select valid test cases (maximal outbreak only)
sel = data.maximal_outbreak.flatten()  # boolean selection mask

# Step 4: compute metrics on selected cases
mrr  = rank_score(true_source_ranks.flatten(), sel)
top5 = top_k_score(true_source_ranks.flatten(), sel, k=5)
\end{lstlisting}
\end{codewalk}

% ============================================================
\section{The Computational Challenge}
\label{sec:computational-challenge}
% ============================================================

We have established the correct approach (MLE via simulation) and the
metrics to evaluate it. Now let us be precise about \emph{why} this
approach cannot scale to real-world networks without clever approximations.

\subsection{Back-of-the-Envelope Calculation}

\begin{examplebox}{Scaling Analysis: How Long Does Brute-Force MLE Take?}
Let us parametrize the cost. Suppose:
\begin{itemize}
    \item $N$ = number of nodes (network size)
    \item $n_{\text{runs}}$ = number of simulation runs per candidate
    \item $t_{\text{sim}}$ = time to run one epidemic simulation
\end{itemize}

Total cost = $N \times n_{\text{runs}} \times t_{\text{sim}}$.

\bigskip
\textbf{Realistic values:}
A modern event-driven SIR simulator (like the C implementation in
\texttt{tsir/}) can simulate one epidemic in roughly $1\,\text{ms}$ for
$N = 1000$ nodes. For accuracy, we need at least $n_{\text{runs}} = 100$
simulations per candidate.

\medskip
\begin{center}
\begin{tabular}{rrrr}
\toprule
\textbf{Network size $N$} & \textbf{Simulations} & \textbf{Time/sim} & \textbf{Total time} \\
\midrule
100     & $100 \times 100 = 10{,}000$       & $0.1\,\text{ms}$ & $1\,\text{second}$    \\
1,000   & $1000 \times 100 = 100{,}000$     & $1\,\text{ms}$   & $100\,\text{seconds}$ \\
10,000  & $10{,}000 \times 100 = 1{,}000{,}000$ & $10\,\text{ms}$ & $\approx 2.8\,\text{hours}$ \\
100,000 & $10^7$                            & $100\,\text{ms}$ & $\approx 115\,\text{days}$ \\
\bottomrule
\end{tabular}
\end{center}

\medskip
For $N = 1000$ nodes, brute-force MLE takes about 100 seconds — feasible
for occasional research use but far too slow for real-time outbreak response.
For $N = 10{,}000$ it is almost 3 hours. For $N = 100{,}000$ (the size of a
small city's contact network), it would take nearly 4 months.
\end{examplebox}

\subsection{Why We Need Faster Methods}

The scaling analysis above makes clear that brute-force MLE is impractical
for any real application. This motivates two alternative strategies, both
of which are developed in this thesis:

\begin{conceptbox}{Two Routes to Scalable Source Detection}
\textbf{1. Classical heuristics (Chapter~5).} Replace the expensive
simulation-based likelihood with a cheap proxy that can be computed in
polynomial time. Examples: the Jordan center computes eccentricities in
$\mathcal{O}(N^2)$ time; rumor centrality runs in $\mathcal{O}(N)$ on trees.
These methods are fast but limited: they rely on strong structural
assumptions (tree topology, SI dynamics) that rarely hold in practice.

\textbf{2. Learning-based methods (Chapters~6--11).} Train a graph neural
network on thousands of simulated epidemics (generated offline). At
inference time, the trained model evaluates each new epidemic observation
in a \emph{single forward pass} of the neural network, requiring only
$\mathcal{O}(|E| + N)$ operations — no simulation needed at test time.
The training cost is paid once; the inference cost is negligible.
\end{conceptbox}

The key insight is that the GNN is \textbf{amortising} the simulation cost:
instead of running 100 simulations per candidate for each new epidemic, we
run millions of simulations \emph{once} during training, and the network
learns to extract the relevant patterns. At test time, a single forward pass
— typically taking milliseconds — replaces the hours of brute-force computation.

\begin{keyinsight}{Why GNNs Are the Right Tool}
Source detection on a graph has a natural structure that GNNs are designed
to exploit:
\begin{enumerate}
    \item The observation $\mathbf{X}$ assigns features to each node
    ($S$, $I$, or $R$) — exactly the input format for a node classification
    GNN.
    \item The posterior $P(\text{source}=v \mid \mathbf{X})$ is a
    probability assigned to each node — exactly the output format for a
    node-level prediction task.
    \item The relevant information is \emph{local}: a node's probability of
    being the source depends primarily on the states of its neighbours and
    the topology of its local neighbourhood, not on distant parts of the
    graph. GNN message passing naturally aggregates this local information.
    \item With temporal networks, edge timestamps provide additional causal
    constraints that temporal GNN architectures can learn to exploit.
\end{enumerate}
\end{keyinsight}

% ============================================================
\section{The Inverse Problem: A Structural View}
\label{sec:structural-view}
% ============================================================

Before closing this chapter, it is worth stepping back to understand the
structure of the problem from a higher level.

\subsection{The Forward and Inverse Directions}

\begin{center}
\begin{tikzpicture}[
  box/.style={draw, rounded corners=4pt, minimum width=3cm,
              minimum height=1cm, font=\small, fill=blue!6, thick,
              align=center},
  arr/.style={-{Stealth[length=3mm]}, thick},
  darr/.style={-{Stealth[length=3mm]}, thick, dashed, gray!70}
]

\node[box] (source) at (0,0)   {Source node $s^*$ \\ \small(unknown)};
\node[box] (process) at (4.5,0) {SIR spreading \\ \small on network $G$};
\node[box] (obs)    at (9,0)   {Observation $\mathbf{X}$ \\ \small(known)};

\draw[arr] (source) -- node[above, font=\footnotesize]{forward / easy} (process);
\draw[arr] (process) -- node[above, font=\footnotesize]{stochastic} (obs);

\draw[darr, bend right=30] (obs) to
  node[below, font=\footnotesize, align=center]{inverse / hard \\ (source detection)} (source);

\end{tikzpicture}
\end{center}

The \textbf{forward direction} (left to right) is tractable: given $s^*$ and
$G$, we simulate the SIR process to generate $\mathbf{X}$. This is what the
TSIR engine does.

The \textbf{inverse direction} (right to left) is what we solve in this
thesis: given $\mathbf{X}$ and $G$, recover $s^*$. The arrow is dashed
because the mapping is not a function — the same $\mathbf{X}$ can arise from
multiple different sources, and the same source can produce many different
$\mathbf{X}$ values. The inverse is a probability distribution, not a
deterministic map.

\subsection{What Makes a Source Detectable?}

The following table summarises the structural conditions that make source
detection easier or harder. Each row represents a variable that we
manipulate in the experimental evaluation of Chapter~6.

\begin{conceptbox}{Structural Determinants of Detection Difficulty}
\begin{center}
\renewcommand{\arraystretch}{1.3}
\begin{tabular}{lll}
\toprule
\textbf{Property} & \textbf{Easier detection} & \textbf{Harder detection} \\
\midrule
Network structure     & Tree-like (few cycles)       & Dense (many cycles)        \\
Degree distribution   & Heterogeneous (scale-free)   & Homogeneous (regular, ER)  \\
Observation time $T$  & Early (small $T$)            & Late (large $T$)           \\
Observation coverage  & Full (all nodes)             & Partial (sensor subset)    \\
Temporal resolution   & High (fine timestamps)       & Low (aggregated/static)    \\
Epidemic size $k/N$   & Intermediate                 & Very small or very large   \\
\bottomrule
\end{tabular}
\end{center}
\end{conceptbox}

Each row of this table is a research question. Our experimental setup in
Chapter~6 varies these conditions systematically to understand when and why
different methods succeed or fail.

% ============================================================
\section{Chapter Summary}
\label{sec:source-detection-summary}
% ============================================================

This chapter introduced the source detection problem from the ground up. We
started with the detective analogy and built toward a rigorous mathematical
framework. The key takeaways are:

\begin{enumerate}
    \item \textbf{Source detection is a Bayesian inverse problem.} The
    posterior $P(\text{source}=v \mid \mathbf{X}) \propto
    P(\mathbf{X} \mid \text{source}=v)$ (under a uniform prior) assigns a
    probability to each candidate node.

    \item \textbf{The likelihood is intractable exactly but approximable by
    simulation.} Running $n_{\text{runs}}$ epidemics from each candidate and
    counting matches (exact or soft via Jaccard similarity) gives an unbiased
    estimate.

    \item \textbf{Real observations are noisy and partial.} Soft matching
    and appropriate handling of partial observations are essential for
    robust methods.

    \item \textbf{Four core challenges.} Partial observation, late
    observation, network heterogeneity, and temporal ambiguity each degrade
    detection performance in distinct ways.

    \item \textbf{Evaluation uses rank score and top-$k$ score.} Both
    metrics are normalised to be network-size-independent and are computed
    only on valid test cases via the \texttt{possible} and
    \texttt{maximal\_outbreak} filters.

    \item \textbf{Brute-force MLE is computationally infeasible} for
    networks beyond a few hundred nodes, motivating both classical heuristics
    (Chapter~5) and GNN-based learning methods (Chapters~6--11).
\end{enumerate}

\begin{keyinsight}{The Thesis in One Sentence}
Source detection is a hard inverse problem over a stochastic spreading
process; Graph Neural Networks trained on simulated epidemic data can learn
to approximate the posterior efficiently in a single forward pass, exploiting
both the graph structure and — in temporal architectures — the causal
constraints imposed by contact timestamps.
\end{keyinsight}
